\chapter{Présentation générale}
\label{chap:presentation}

\sectionsansnum{Introduction}

Ce premier chapitre situe le projet dans son double contexte : celui de l'organisme
qui l'a accueilli et celui de la médecine du sommeil. Nous présentons le cadre du
stage, l'organisme d'accueil et son activité, puis les notions médicales
indispensables à la compréhension du mémoire. De ce contexte émerge une problématique
que nous confrontons aux solutions existantes, avant d'en déduire le périmètre de la
solution proposée et la méthodologie retenue pour la construire.

\section{Cadre général du stage}

Ce projet de fin d'études s'est déroulé du \textbf{26 janvier au 26 juillet 2026},
soit six mois, au sein du \textbf{Centre National de l'Informatique} à Tunis. Il a été
encadré côté entreprise par M. Abdelouaheb Dahmeni et M\textsuperscript{me} Hajer
Dahmeni, et côté académique par M\textsuperscript{me} Ines Chouchane.

Le projet consiste à concevoir et réaliser \hypnoria{}, une plateforme web d'aide à
l'interprétation des enregistrements polysomnographiques. Il couvre l'ensemble du
cycle de développement : constitution du corpus de données, entraînement et évaluation
de modèles d'apprentissage, conception et développement de l'application, puis
conteneurisation et automatisation du déploiement.

\section{Présentation de l'organisme d'accueil}

\subsection{Le Centre National de l'Informatique}

Le Centre National de l'Informatique (CNI) est un établissement public à caractère non
administratif, doté de la personnalité civile et de l'autonomie financière. Il a été
créé par la loi n\textsuperscript{o} 75-83 du 30 décembre 1975, dont les articles 35 à
42 fixent son statut et ses attributions. Placé sous la tutelle du ministère chargé des
technologies de la communication, il est certifié ISO 9001 version 2015.

Le CNI constitue le principal appui technique des structures publiques tunisiennes pour
la conception, le déploiement et l'exploitation de leurs systèmes d'information. Il
intervient aussi bien en maîtrise d'ouvrage déléguée qu'en réalisation directe, et
opère une infrastructure d'hébergement mutualisée au service de l'administration.

Le CNI a accompagné toutes les vagues d'informatisation de l'administration tunisienne,
depuis les traitements centralisés des années 1970 jusqu'aux services numériques
contemporains. La figure~\ref{fig:cni-timeline} en retrace les principaux jalons.

\begin{figure}[H]
\centering
\begin{tikzpicture}[x=1cm,y=1cm,
    jalon/.style   = {align=center,font=\scriptsize,text width=2.6cm},
    millesime/.style = {font=\footnotesize\bfseries,text=hypnoblue},
  ]

  % Axe temporel
  \draw[line width=1.1pt,hypnoblue,-{Latex[length=3mm]}] (-0.4,0) -- (15,0);

  % Jalons portés au-dessus de l'axe
  \foreach \px/\yr/\txt in {
      0/1975/{Création du CNI\\par la loi 75-83},
      6/1992/{Élargissement de\\l'offre de formation},
      12/2007/{Réseau national,\\e-gouvernement, backup}}
  {
    \draw[line width=1pt,hypnoblue] (\px,-0.18) -- (\px,0.18);
    \fill[hypnored] (\px,0) circle (2.4pt);
    \node[jalon,anchor=south]     at (\px, 0.62) {\txt};
    \node[millesime,anchor=north] at (\px,-0.24) {\yr};
  }

  % Jalons portés sous l'axe
  \foreach \px/\yr/\txt in {
      3/1978/{Prise en charge de la\\formation en informatique},
      9/1997/{Reclassement en\\établissement public},
      14.4/{Auj.}/{Pôle d'excellence\\des TIC publiques}}
  {
    \draw[line width=1pt,hypnoblue] (\px,-0.18) -- (\px,0.18);
    \fill[hypnored] (\px,0) circle (2.4pt);
    \node[jalon,anchor=north]     at (\px,-0.62) {\txt};
    \node[millesime,anchor=south] at (\px, 0.24) {\yr};
  }

\end{tikzpicture}
\caption[Chronologie du Centre National de l'Informatique]%
        {Principaux jalons de l'évolution du Centre National de l'Informatique.}
\label{fig:cni-timeline}
\end{figure}

\subsection{Services proposés par le CNI}

Les missions du CNI couvrent l'ensemble du cycle de vie d'un système d'information
public :

\begin{itemize}[itemsep=2pt,topsep=4pt]
  \item \textbf{Maîtrise d'ouvrage déléguée} --- pilotage de projets pour le compte
        des administrations bénéficiaires.
  \item \textbf{Études et conseil} --- audits, études d'opportunité, schémas
        directeurs, rédaction de cahiers des charges, conseil en réseaux et en sécurité
        des systèmes d'information.
  \item \textbf{Développement} --- conception, réalisation et maintenance applicative,
        accompagnée de la formation des utilisateurs aux systèmes livrés.
  \item \textbf{Hébergement} --- exploitation de serveurs, d'applications et de données
        pour le compte des organismes publics.
  \item \textbf{Continuité d'activité} --- solutions de sauvegarde, plans de reprise
        après sinistre et plans de continuité.
  \item \textbf{Formation et certification} --- formations continues et certifiantes,
        séminaires techniques.
  \item \textbf{Déploiement et assistance} --- mise en production des applications,
        support sur site ou à distance.
\end{itemize}

\subsection{Intégration du projet au sein du CNI}

Le projet \hypnoria{} s'inscrit dans les axes « Études et conseil » et
« Développement » du CNI. Il relève d'une démarche exploratoire : appliquer les
techniques d'apprentissage automatique récentes à l'aide au diagnostic médical, domaine
où l'administration tunisienne dispose de gisements de données mais de peu d'outils
d'exploitation.

Trois caractéristiques du contexte CNI ont orienté les choix techniques du projet :

\begin{itemize}[itemsep=2pt,topsep=4pt]
  \item \textbf{La souveraineté des données.} Les données de santé sont sensibles.
        L'architecture retenue sépare strictement les métadonnées cliniques, stockées
        dans une base relationnelle maîtrisée, des fichiers volumineux, confiés à un
        stockage objet dont l'accès n'est possible que par des liens signés à durée
        limitée.
  \item \textbf{La reproductibilité du déploiement.} Le CNI exploite des systèmes sur
        une longue durée. L'ensemble de la plateforme est donc conteneurisé et
        démarrable en une seule commande.
  \item \textbf{La traçabilité.} Toute action sensible est journalisée avec son auteur,
        son horodatage et son adresse d'origine, et le journal est exportable en vue
        d'un contrôle de conformité.
\end{itemize}

\section{Présentation du projet}

\subsection{Le sommeil et sa structure}

Le sommeil n'est pas un état homogène de repos mais une succession organisée d'états
physiologiques distincts, qui se répètent au cours de la nuit selon des cycles d'une
durée moyenne de 90 minutes. On distingue le sommeil sans mouvements oculaires rapides,
ou \emph{NREM}, subdivisé en stades de profondeur croissante, et le sommeil paradoxal,
ou \emph{REM}, caractérisé par une activité cérébrale proche de l'éveil associée à une
atonie musculaire complète.

La classification de référence est celle de l'\emph{American Academy of Sleep
Medicine}~\cite{aasm2020}, qui a succédé en 2007 aux règles de Rechtschaffen et
Kales~\cite{rechtschaffen1968}. Elle définit cinq états, décrits au
tableau~\ref{tab:stades}, et impose que le scorage se fasse par fenêtres, ou
\emph{époques}, de 30 secondes.

\begin{table}[H]
\centering

\small
\begin{tabular}{@{}lL{4.6cm}L{5.4cm}c@{}}
\toprule
\textbf{Stade} & \textbf{Signature électrophysiologique} & \textbf{Éléments
caractéristiques} & \textbf{Part typique} \\
\midrule
W (éveil) & Rythme alpha postérieur, activité bêta, tonus musculaire élevé &
Mouvements oculaires volontaires et clignements & --- \\
\addlinespace[2pt]
N1 & Ralentissement du rythme de fond, ondes thêta & Mouvements oculaires lents,
transition instable & \SI{5}{\percent} \\
\addlinespace[2pt]
N2 & Activité thêta dominante & Fuseaux du sommeil et complexes K & \SI{45}{\percent}
à \SI{55}{\percent} \\
\addlinespace[2pt]
N3 & Ondes lentes de grande amplitude, sur plus de \SI{20}{\percent} de l'époque &
Sommeil profond, seuil d'éveil élevé & \SI{15}{\percent} à \SI{25}{\percent} \\
\addlinespace[2pt]
REM & Rythme rapide et désynchronisé, proche de l'éveil & Mouvements oculaires
rapides, atonie musculaire & \SI{20}{\percent} à \SI{25}{\percent} \\
\bottomrule
\end{tabular}
\caption{Les cinq états de vigilance définis par le référentiel AASM}
\label{tab:stades}
\end{table}

Deux points conditionnent la difficulté de la tâche d'automatisation. D'une part, la
distribution des stades est fortement déséquilibrée : le N2 représente à lui seul
environ la moitié d'une nuit, tandis que le N1 en occupe à peine un vingtième. D'autre
part, le stade N1 est un état de transition sans marqueur graphique propre : il se
définit par la disparition du rythme alpha et l'absence des fuseaux du N2, c'est-à-dire
négativement. C'est la classe sur laquelle les experts humains eux-mêmes s'accordent le
moins.

La succession des stades attribués aux époques successives constitue
l'\emph{hypnogramme} (figure~\ref{fig:hypnogramme}), document de référence du médecin
du sommeil : il condense environ mille décisions élémentaires en une seule courbe.

\begin{figure}[H]
\centering
\begin{tikzpicture}[x=1.02cm,y=1cm,
    annot/.style = {font=\scriptsize,align=center},
    fleche/.style= {-{Latex[length=2mm]},line width=0.55pt},
  ]

  % ── Bandes signalant les épisodes de sommeil paradoxal ────────────────────
  \begin{scope}[on background layer]
    \foreach \a/\b in {2.25/2.625, 4.5/5.025, 7.125/7.875, 9.0/9.9, 10.95/12.0}
      \fill[stREM!14] (\a,-3.95) rectangle (\b,0.30);
  \end{scope}

  % ── Lignes de niveau ──────────────────────────────────────────────────────
  \foreach \y in {0,-0.9,-1.8,-2.7,-3.6}
    \draw[hypnogray!22,line width=0.4pt] (0,\y) -- (12.3,\y);

  % ── Étiquettes des stades ─────────────────────────────────────────────────
  \node[left,font=\scriptsize\bfseries,text=stWake] at (-0.12,0)    {W};
  \node[left,font=\scriptsize\bfseries,text=stREM]  at (-0.12,-0.9) {REM};
  \node[left,font=\scriptsize\bfseries,text=stN1]   at (-0.12,-1.8) {N1};
  \node[left,font=\scriptsize\bfseries,text=stN2]   at (-0.12,-2.7) {N2};
  \node[left,font=\scriptsize\bfseries,text=stN3]   at (-0.12,-3.6) {N3};

  % ── Tracé ─────────────────────────────────────────────────────────────────
  \draw[line width=1.15pt,hypnoblue]
    (0,0) -- (0.45,0) -- (0.45,-1.8) -- (0.675,-1.8) -- (0.675,-2.7) --
    (1.275,-2.7) -- (1.275,-3.6) -- (1.95,-3.6) -- (1.95,-2.7) --
    (2.25,-2.7) -- (2.25,-0.9) -- (2.625,-0.9) -- (2.625,-1.8) --
    (2.775,-1.8) -- (2.775,-2.7) -- (3.45,-2.7) -- (3.45,-3.6) --
    (4.2,-3.6) -- (4.2,-2.7) -- (4.5,-2.7) -- (4.5,-0.9) --
    (5.025,-0.9) -- (5.025,0) -- (5.175,0) -- (5.175,-1.8) --
    (5.325,-1.8) -- (5.325,-2.7) -- (6.3,-2.7) -- (6.3,-3.6) --
    (6.75,-3.6) -- (6.75,-2.7) -- (7.125,-2.7) -- (7.125,-0.9) --
    (7.875,-0.9) -- (7.875,0) -- (8.025,0) -- (8.025,-1.8) --
    (8.175,-1.8) -- (8.175,-2.7) -- (9.0,-2.7) -- (9.0,-0.9) --
    (9.9,-0.9) -- (9.9,0) -- (10.125,0) -- (10.125,-1.8) --
    (10.35,-1.8) -- (10.35,-2.7) -- (10.95,-2.7) -- (10.95,-0.9) --
    (12.0,-0.9) -- (12.0,0) -- (12.3,0);

  % ── Axe des temps ─────────────────────────────────────────────────────────
  \draw[hypnogray,line width=0.6pt] (0,-4.05) -- (12.3,-4.05);
  \foreach \h in {0,1,2,3,4,5,6,7,8} {
    \pgfmathsetmacro{\xx}{\h*1.5}
    \draw[hypnogray,line width=0.5pt] (\xx,-4.05) -- (\xx,-4.22);
    \node[below,font=\scriptsize,text=hypnogray,inner sep=1.5pt] at (\xx,-4.22) {\h\,h};
  }

  % ── Annotation : sommeil profond ──────────────────────────────────────────
  \node[annot,text=hypnored,text width=3.9cm,anchor=north] at (2.6,-4.85)
       {sommeil profond concentré\\dans le premier tiers de la nuit};
  \draw[fleche,hypnored] (1.75,-4.80) -- (1.62,-3.72);

  % ── Annotation : sommeil paradoxal ────────────────────────────────────────
  \node[annot,text=stREM,text width=4.3cm,anchor=south] at (9.7,1.05)
       {épisodes de sommeil paradoxal\\de plus en plus longs};
  \draw[fleche,stREM] (10.75,1.00) -- (11.45,-0.78);

\end{tikzpicture}
\vspace{0.35cm}
\caption[Hypnogramme d'une nuit physiologique]%
        {Hypnogramme schématique d'une nuit physiologique de huit heures. Les bandes
         vertes signalent les épisodes de sommeil paradoxal.}
\label{fig:hypnogramme}
\end{figure}

\subsection{La polysomnographie}

L'hypnogramme n'est pas mesuré, il est \emph{interprété}. Il découle d'un examen appelé
polysomnographie, qui enregistre simultanément plusieurs signaux physiologiques pendant
une nuit complète. Le tableau~\ref{tab:derivations} décrit les dérivations minimales
requises pour établir un scorage, qui sont précisément celles exploitées dans ce
travail.

\begin{table}[H]
\centering

\small
\begin{tabular}{@{}llL{7.2cm}@{}}
\toprule
\textbf{Signal} & \textbf{Nombre} & \textbf{Apport au scorage} \\
\midrule
Électroencéphalogramme (EEG) & 2 & Rythme de fond, fuseaux, complexes K et ondes
lentes : c'est le signal qui porte la décision \\
\addlinespace[2pt]
Électro-oculogramme (EOG) & 2 & Distingue les mouvements oculaires lents de
l'endormissement des mouvements rapides du sommeil paradoxal \\
\addlinespace[2pt]
Électromyogramme (EMG) mentonnier & 1 & Le tonus musculaire sépare le sommeil
paradoxal du sommeil lent : son effondrement signe le REM \\
\bottomrule
\end{tabular}
\caption{Dérivations minimales d'un montage polysomnographique de scorage}
\label{tab:derivations}
\end{table}

Un montage complet comporte en pratique d'autres capteurs --- flux nasal, sangles
thoraciques, oxymétrie de pouls, électrocardiogramme --- destinés à la caractérisation
des événements respiratoires. Le scorage des stades, lui, ne requiert que les cinq
dérivations ci-dessus. Ce point est structurant : il rend possible un système
compatible avec des montages allégés, et nous évaluerons d'ailleurs une configuration
réduite à deux dérivations.

\subsection{Le syndrome d'apnées obstructives du sommeil}

Le syndrome d'apnées obstructives du sommeil (SAOS) résulte d'obstructions répétées des
voies aériennes supérieures pendant le sommeil. Chaque obstruction provoque une
désaturation en oxygène et un micro-éveil, invisible pour le patient mais destructeur
pour l'architecture du sommeil.

Sa prévalence est considérable : une analyse publiée en 2019 estime à environ
936 millions le nombre d'adultes de 30 à 69 ans présentant un SAOS léger à sévère dans
le monde, dont près de 425 millions sous une forme modérée à
sévère~\cite{benjafield2019}. Non traité, le syndrome est associé à une surmorbidité
cardiovasculaire documentée~\cite{marin2005}, et il reste largement
sous-diagnostiqué~\cite{young1997}.

La sévérité s'exprime par l'index d'apnées-hypopnées, nombre moyen d'événements
respiratoires par heure de sommeil. Le référentiel AASM~\cite{kapur2017} en définit
quatre classes, reprises au tableau~\ref{tab:iah}, qui constituent les classes cibles
du second modèle développé dans ce mémoire.

\begin{table}[H]
\centering

\small
\begin{tabular}{@{}lcL{7.4cm}@{}}
\toprule
\textbf{Classe} & \textbf{IAH (évén./h)} & \textbf{Conduite habituelle} \\
\midrule
Normal   & $< 5$          & Pas de traitement spécifique \\
Léger    & $5$ à $< 15$   & Mesures hygiéno-diététiques, surveillance \\
Modéré   & $15$ à $< 30$  & Traitement discuté selon la symptomatologie \\
Sévère   & $\geq 30$      & Traitement par pression positive continue indiqué \\
\bottomrule
\end{tabular}
\caption{Classes de sévérité du SAOS selon l'index d'apnées-hypopnées}
\label{tab:iah}
\end{table}

Le point essentiel pour la suite est le suivant : le SAOS ne se contente pas de
provoquer des événements respiratoires, il \emph{déforme l'architecture du sommeil}.
Les micro-éveils répétés fragmentent la nuit, suppriment le sommeil profond et
raccourcissent les épisodes paradoxaux. Cette empreinte est mesurable sur le seul
hypnogramme, indépendamment des capteurs respiratoires --- c'est l'hypothèse sur
laquelle repose le second volet du projet.

\subsection{Objectif général du projet}

\hypnoria{} est une plateforme web d'aide à l'interprétation polysomnographique
destinée aux médecins du sommeil. Elle poursuit quatre objectifs :

\begin{enumerate}[itemsep=3pt,topsep=4pt]
  \item \textbf{Automatiser le scorage des stades} à partir d'un fichier
        d'enregistrement brut au format EDF, avec une exactitude comparable à l'accord
        inter-expert.
  \item \textbf{Prédire la sévérité du SAOS} à partir de l'architecture du sommeil et
        de marqueurs cliniques, sans exiger l'annotation manuelle des événements
        respiratoires.
  \item \textbf{Expliquer chaque prédiction} en exposant les facteurs qui l'ont
        déterminée, afin que le praticien puisse la valider ou la contester.
  \item \textbf{Fournir la chaîne applicative complète} : dossier patient, archivage
        sécurisé, annotation, compte rendu exportable, collaboration et administration
        multi-établissements.
\end{enumerate}

\subsection{Problématique}

Trois constats fondent ce projet.

\subsubsection{Le coût du scorage manuel}

Une nuit d'enregistrement de huit heures produit environ un millier d'époques de
30 secondes. Chacune doit être examinée, comparée à ses voisines et classée par un
technicien qualifié. L'opération représente plusieurs heures de travail expert par
patient, auxquelles s'ajoutent l'analyse des événements respiratoires et la rédaction du
compte rendu. Dans un service saturé, le facteur limitant n'est plus le nombre de lits
d'enregistrement mais le temps d'expertise disponible pour interpréter ce qu'ils
produisent.

\subsubsection{La variabilité inter-scoreur}

Le scorage manuel n'est pas seulement coûteux, il est également imparfaitement
reproductible. Le programme de fiabilité inter-scoreurs de l'AASM mesure un accord
global de l'ordre de \SI{83}{\percent}~\cite{rosenberg2013}, avec des écarts marqués
selon le stade~\cite{danker2009}.

Deux conséquences en découlent. La première est méthodologique : les annotations de
référence ne sont pas une vérité absolue mais l'opinion d'un scoreur, et l'accord
inter-expert constitue donc un plafond de performance réaliste. La seconde est
clinique : elle légitime l'automatisation, car un modèle présente au moins la propriété
d'être parfaitement reproductible.

\subsubsection{L'opacité des modèles d'apprentissage profond}

Depuis 2017, les réseaux de neurones profonds atteignent sur les bases publiques des
performances comparables à celles d'un scoreur humain~\cite{supratak2017,perslev2021}.
Le verrou n'est donc plus principalement la précision : il est l'adoption.

Un modèle d'apprentissage profond fonctionne comme une boîte noire : il produit une
classe sans énoncer ce qui l'a déterminée. Or le praticien engage sa responsabilité sur
le diagnostic qu'il signe, et ne peut ni valider ni contester une décision dont il
ignore les motifs. À ce frein pratique s'ajoute un cadre réglementaire européen qui
reconnaît un droit à l'explication pour les décisions automatisées~\cite{goodman2017}.

\medskip
\noindent
Ces trois constats convergent : la médecine du sommeil a besoin d'un système qui soit à
la fois \textbf{rapide}, pour absorber le volume d'examens, \textbf{reproductible},
pour lever la variabilité inter-scoreur, et \textbf{transparent}, pour que le praticien
puisse s'en porter garant.

\section{Étude préalable}

\subsection{Étude de l'existant}

\subsubsection{Les solutions commerciales}

Le marché de l'analyse polysomnographique est dominé par des logiciels propriétaires
adossés à du matériel d'acquisition : \emph{Noxturnal} de Nox Medical, \emph{Sleepware
G3} de Philips Respironics, \emph{ProFusion PSG} de Compumedics ou encore
\emph{SleepWorks} de Natus. Une seconde génération d'acteurs, dont \emph{EnsoSleep}
d'EnsoData, propose un scorage automatique fondé sur l'apprentissage profond.

En Tunisie, les laboratoires du sommeil s'appuient sur ces mêmes suites, acquises avec
le matériel d'enregistrement. Le coût de licence et l'adhérence au matériel du
constructeur en limitent la diffusion, et aucune solution nationale d'analyse
automatique n'est à notre connaissance disponible.

\subsubsection{Les travaux académiques}

La littérature scientifique a produit une succession d'architectures de plus en plus
performantes, résumées au tableau~\ref{tab:etat-art-staging}.
\emph{DeepSleepNet}~\cite{supratak2017} combine une extraction convolutive
multi-échelle et une modélisation séquentielle récurrente sur un unique canal.
\emph{TinySleepNet}~\cite{supratak2020} en propose une version allégée.
\emph{U-Sleep}~\cite{perslev2021} démontre une robustesse remarquable au montage et à
la cohorte. \emph{SleepTransformer}~\cite{phan2022} introduit un mécanisme
d'auto-attention ainsi qu'une forme d'interprétabilité par les poids d'attention.
\emph{XSleepNet}~\cite{phan2021} exploite conjointement la représentation temporelle
brute et la représentation temps-fréquence.

\begin{table}[H]
\centering

\small
\begin{tabular}{@{}L{3.3cm}cL{4.4cm}L{2.2cm}c@{}}
\toprule
\textbf{Travail} & \textbf{Année} & \textbf{Approche} & \textbf{Évaluation} &
\textbf{Exactitude} \\
\midrule
Fraiwan \emph{et al.}~\cite{fraiwan2012} & 2012 & Temps-fréquence + forêt aléatoire &
jeu propriétaire & $\approx$ \SI{83}{\percent} \\
\addlinespace[2pt]
Tsinalis \emph{et al.}~\cite{tsinalis2016} & 2016 & CNN sur EEG brut & Sleep-EDF &
$\approx$ \SI{82}{\percent} \\
\addlinespace[2pt]
DeepSleepNet~\cite{supratak2017} & 2017 & CNN multi-échelle + BiLSTM & Sleep-EDF &
$\approx$ \SI{82}{\percent} \\
\addlinespace[2pt]
Chambon \emph{et al.}~\cite{chambon2018} & 2018 & CNN multivarié (EEG, EOG, EMG) &
MASS & $\approx$ \SI{83}{\percent} \\
\addlinespace[2pt]
SleepEEGNet~\cite{mousavi2019} & 2019 & CNN + séquence à séquence & Sleep-EDF &
$\approx$ \SI{84}{\percent} \\
\addlinespace[2pt]
TinySleepNet~\cite{supratak2020} & 2020 & CNN + LSTM allégés & Sleep-EDF &
$\approx$ \SI{85}{\percent} \\
\addlinespace[2pt]
AttnSleep~\cite{eldele2021} & 2021 & Multi-résolution + auto-attention & Sleep-EDF &
$\approx$ \SI{84}{\percent} \\
\addlinespace[2pt]
U-Sleep~\cite{perslev2021} & 2021 & Entièrement convolutif & 15 cohortes &
F1 macro $\approx$ 0,79 \\
\addlinespace[2pt]
SleepTransformer~\cite{phan2022} & 2022 & Transformer hiérarchique & SHHS &
$\approx$ \SI{88}{\percent} \\
\addlinespace[2pt]
XSleepNet~\cite{phan2021} & 2022 & Fusion brut + temps-fréquence & SHHS &
$\approx$ \SI{89}{\percent} \\
\bottomrule
\end{tabular}
\caption{Principaux travaux de classification automatique des stades du sommeil}
\label{tab:etat-art-staging}
\end{table}

\subsection{Critique de l'existant}

Les suites commerciales sont matures et validées cliniquement, mais présentent quatre
limites au regard de notre problématique. L'\textbf{adhérence matérielle} contraint les
laboratoires disposant d'un parc hétérogène. Le \textbf{coût de licence}, calibré pour
des structures hospitalières occidentales, reste hors de portée de nombreux centres.
L'\textbf{installation locale} est peu compatible avec le travail à distance et la
collaboration entre praticiens de sites différents. Enfin, lorsqu'un scorage
automatique est proposé, il l'est \textbf{sans explication} des facteurs qui l'ont
produit.

Les travaux académiques, de leur côté, visent la performance de classification évaluée
sur des jeux publics et s'arrêtent au modèle. Aucun ne livre la chaîne applicative qui
permettrait à un praticien d'en faire usage : gestion des patients, archivage des
examens, production d'un compte rendu, collaboration.

Le tableau~\ref{tab:existant} confronte ces deux familles aux besoins identifiés.

\begin{table}[H]
\centering

\small
\begin{tabular}{@{}L{4.5cm}ccc@{}}
\toprule
\textbf{Critère} & \textbf{Suites commerciales} & \textbf{Travaux académiques} &
\textbf{\hypnoria{}} \\
\midrule
Scorage automatique            & partiel  & oui  & oui \\
Indépendance au matériel       & faible   & --   & oui \\
Explicabilité des décisions    & non      & rare & oui \\
Prédiction de la sévérité SAOS & indirecte& non  & oui \\
Application web multi-utilisateurs & non  & non  & oui \\
Dossier patient et historique  & oui      & non  & oui \\
Collaboration entre praticiens & non      & non  & oui \\
Traçabilité et audit           & partielle& non  & oui \\
Coût d'acquisition             & élevé    & nul  & nul \\
\bottomrule
\end{tabular}
\caption{Positionnement de \hypnoria{} face aux solutions existantes}
\label{tab:existant}
\end{table}

\subsection{Solution proposée}

Il ressort de cette analyse un espace non couvert, à l'intersection des deux familles :
une plateforme applicative complète, indépendante du matériel d'acquisition, qui intègre
des modèles de l'état de l'art \emph{et} les explique. C'est cet espace que le projet
occupe.

Le système s'articule autour de deux étages d'apprentissage successifs et d'une
plateforme qui les met à disposition. Le premier étage assure la \textbf{classification
des stades} : trois familles d'architectures --- convolutive, récurrente et
attentionnelle --- sont entraînées séparément puis combinées par un méta-apprenant. Deux
configurations de montage et deux granularités de sortie sont proposées, afin de couvrir
aussi bien l'enregistrement de laboratoire que le dispositif ambulatoire. Le second
étage assure la \textbf{prédiction de la sévérité} : il extrait de l'hypnogramme
produit une vingtaine de marqueurs d'architecture du sommeil, les complète de données
cliniques et oxymétriques, et prédit la classe de sévérité en accompagnant chaque
prédiction de sa décomposition en contributions par variable.

\section{Méthodologie de travail}

\subsection{Choix du cadre Scrum}

Le projet présentait deux sources d'incertitude difficilement compatibles avec une
planification linéaire. La première tenait aux données : l'accès aux cohortes cliniques
est soumis à une autorisation dont le délai d'instruction n'était pas connu à l'avance.
La seconde tenait aux modèles : rien ne garantissait \emph{a priori} qu'une
architecture donnée atteindrait le niveau de performance visé, ni que l'empilement
apporterait un gain.

Une démarche séquentielle aurait figé des décisions avant de disposer des éléments
permettant de les prendre. Nous avons donc retenu le cadre Scrum~\cite{schwaber2020},
dont le principe est d'avancer par incréments courts et d'ajuster la suite à ce que
chaque incrément a révélé.

\subsection{Les rôles dans la méthodologie Scrum}

Scrum définit trois rôles. Le \emph{Product Owner} porte la vision du produit, arbitre
les priorités du backlog et valide la valeur livrée. Le \emph{Scrum Master} garantit le
respect du cadre, anime les cérémonies et lève les obstacles. L'\emph{équipe de
développement} réalise l'incrément.

Le projet étant mené par une seule personne, ces rôles ont été redistribués sans être
dénaturés : la fonction de Product Owner a été portée par l'encadrement professionnel du
CNI, celle de Scrum Master par l'encadrement académique lors des points de suivi, et
celle d'équipe de développement par l'étudiant.

\subsection{Adaptation au projet}

La durée d'itération a été portée à quatre semaines plutôt que deux. Ce choix est
directement lié à la nature du travail : l'entraînement d'un modèle profond sur près de
deux cent mille exemples se compte en heures de calcul, et douze variantes devaient être
produites. Une itération de deux semaines n'aurait pas permis de livrer un incrément
évaluable pendant la phase de modélisation.

Trois artefacts ont été tenus à jour : le \emph{Product Backlog}, liste priorisée des
fonctionnalités attendues ; le \emph{Sprint Backlog}, sous-ensemble sélectionné pour
l'itération en cours ; et l'\emph{incrément}, version fonctionnelle livrée en fin de
sprint. Les cérémonies ont été condensées : une réunion de planification en début de
sprint, un point hebdomadaire avec l'encadrement en substitution de la mêlée
quotidienne, une revue de sprint sous forme de démonstration et une rétrospective
écrite.

\begin{figure}[H]
\centering
\begin{tikzpicture}[
    node distance=0pt,
    box/.style   = {rectangle,rounded corners=2pt,draw=hypnoblue!55,fill=hypnolight,
                    align=center,font=\scriptsize,minimum height=1.05cm,text width=2.15cm},
    dark/.style  = {rectangle,rounded corners=2pt,draw=hypnored!70,fill=hypnored!8,
                    align=center,font=\scriptsize,minimum height=1.05cm,text width=2.15cm},
    fl/.style    = {-{Latex[length=2.2mm]},hypnoblue!70,line width=0.8pt},
  ]

  \node[box] (backlog) at (0,0)      {\textbf{Product\\ Backlog}\\[1pt] 60 user stories};
  \node[box] (planif)  at (3.3,0)    {\textbf{Planification}\\[1pt] sélection et
                                      estimation};
  \node[dark] (sprint) at (6.6,0)    {\textbf{Sprint}\\[1pt] 4 semaines};
  \node[box] (revue)   at (9.9,0)    {\textbf{Revue}\\[1pt] démonstration à
                                      l'encadrement};
  \node[box] (retro)   at (13.0,0)   {\textbf{Rétrospective}\\[1pt] ajustement du
                                      sprint suivant};

  \draw[fl] (backlog) -- (planif);
  \draw[fl] (planif)  -- (sprint);
  \draw[fl] (sprint)  -- (revue);
  \draw[fl] (revue)   -- (retro);

  % Boucle de retour
  \draw[fl] (retro.south) -- ++(0,-0.85) -- node[below,font=\scriptsize\itshape,
        pos=0.5,hypnogray] {réalimentation du backlog à chaque itération}
        ++(-13.0,0) -- (backlog.south);

  % Incrément
  \node[font=\scriptsize\itshape,hypnogray,align=center] at (6.6,1.15)
       {incrément fonctionnel livré};
  \draw[fl] (sprint.north) -- (6.6,0.95);

\end{tikzpicture}
\vspace{0.6cm}
\caption[Cycle Scrum adapté au projet]%
        {Cycle Scrum adapté à un projet individuel de six mois.}
\label{fig:scrum}
\end{figure}

\sectionsansnum{Conclusion}

Ce chapitre a posé le cadre du projet. Nous avons présenté le Centre National de
l'Informatique et montré comment trois de ses contraintes structurantes ont orienté les
choix techniques. Nous avons ensuite introduit les notions médicales nécessaires : les
cinq stades du référentiel AASM, l'hypnogramme, le montage polysomnographique et les
classes de sévérité du SAOS.

De ce contexte se dégage une problématique à trois volets --- le coût du scorage manuel,
sa variabilité entre experts et l'opacité des modèles qui prétendent l'automatiser.
L'étude de l'existant a montré que les suites commerciales et les travaux académiques
traitent chacun une partie du problème sans le couvrir entièrement, laissant ouvert
l'espace qu'occupe \hypnoria{}.

Le chapitre suivant traduit ces intentions en une spécification vérifiable : acteurs,
besoins fonctionnels et non fonctionnels, backlog du produit, modèle de données et
architecture globale.
